\chapter{Duality}
\label{chap:duality}

Duality is one of the most efficient organizing principles in exactly solvable
models.  A duality rewrites the same partition function, Hamiltonian, operator
algebra, or low-energy field theory in variables adapted to a different physical
regime.  In the Ising family, the essential effect is the exchange
\begin{equation}
\label{eq:duality-weak-strong-schematic}
\begin{aligned}
    \text{ordered variables}
    &\longleftrightarrow
    \text{domain-wall or disorder variables},\\
    \text{strong coupling}
    &\longleftrightarrow
    \text{weak coupling}.
\end{aligned}
\end{equation}
Thus duality is not merely a computational trick.  It explains why an apparently
nonlocal description can become local after changing variables, why the
transverse-field Ising critical point is located at the self-dual coupling, and
why the two phases of the Ising chain are naturally described as order and
disorder phases.

This chapter develops duality at three levels:
\begin{enumerate}[label=(\roman*)]
    \item Kramers--Wannier duality of the two-dimensional classical Ising model;
    \item the operator Kramers--Wannier duality of the one-dimensional TFIM; and
    \item the precise logical status of self-duality as a criterion for
    criticality.
\end{enumerate}
The conventions are matched to the TFIM Hamiltonian used in these notes,
\begin{equation}
\label{eq:duality-chapter-tfim}
    H_{\rm TFIM}
    =
    -J\sum_j \sigma_j^x\sigma_{j+1}^x
    -h\sum_j \sigma_j^z,
    \qquad
    J,h\geq0 .
\end{equation}
The ordered phase is therefore the phase with long-range order in
$\sigma^x$, while the transverse-field or disordered phase is polarized by
$\sigma^z$.

\section{Classical Kramers--Wannier duality}
\label{sec:classical-kramers-wannier-duality}

Consider the anisotropic two-dimensional classical Ising model on a square
lattice,
\begin{equation}
\label{eq:duality-2d-classical-ising}
    Z(K_x,K_y)
    =
    \sum_{s_{r}=\pm1}
    \exp\left[
        K_x\sum_{\langle r r'\rangle_x}s_rs_{r'}
        +
        K_y\sum_{\langle r r'\rangle_y}s_rs_{r'}
    \right],
\end{equation}
where $K_x$ and $K_y$ are dimensionless couplings along the two lattice
directions.  The high-temperature expansion uses
\begin{equation}
\label{eq:duality-high-temp-bond}
    \ee^{K s_rs_{r'}}
    =
    \cosh K\left(1+s_rs_{r'}\tanh K\right).
\end{equation}
After expanding all bonds and summing over the Ising spins, only configurations
in which an even number of occupied bonds touches each site survive.  These are
closed loops on the original lattice.  Schematically,
\begin{equation}
\label{eq:duality-high-temp-loops}
    Z(K_x,K_y)
    =
    2^N(\cosh K_x)^{E_x}(\cosh K_y)^{E_y}
    \sum_{\mathcal{C}\,\text{closed}}
    (\tanh K_x)^{n_x(\mathcal{C})}
    (\tanh K_y)^{n_y(\mathcal{C})}.
\end{equation}
Here $N$ is the number of spins, $E_x,E_y$ are the numbers of bonds in the two
directions, and $n_x(\mathcal{C}),n_y(\mathcal{C})$ count occupied loop edges.

The low-temperature expansion instead starts from the two ferromagnetic ground
states.  A configuration is specified by domain walls on the dual lattice.  A
domain wall segment crossing a bond of coupling $K$ costs a Boltzmann factor
$\ee^{-2K}$.  Matching the loop weights in the high-temperature expansion to the
domain-wall weights of a dual Ising model gives the dual couplings
\begin{equation}
\label{eq:duality-classical-dual-couplings-exp}
    \ee^{-2K_x^*}=\tanh K_y,
    \qquad
    \ee^{-2K_y^*}=\tanh K_x .
\end{equation}
Equivalently,
\begin{equation}
\label{eq:duality-classical-dual-couplings-sinh}
    \sinh(2K_x^*)\sinh(2K_y)=1,
    \qquad
    \sinh(2K_y^*)\sinh(2K_x)=1 .
\end{equation}
The exchange of directions in \cref{eq:duality-classical-dual-couplings-exp}
comes from the fact that a horizontal bond of the dual lattice crosses a vertical
bond of the original lattice, and conversely.

Up to a non-singular prefactor depending on the couplings and the boundary
conditions, the partition function of the original model equals the partition
function of the dual model,
\begin{equation}
\label{eq:duality-classical-partition-map}
    Z(K_x,K_y)
    \propto
    Z(K_x^*,K_y^*) .
\end{equation}
The proportionality factor affects the analytic part of the free energy but not
the location of a nonanalytic singularity in the thermodynamic limit.

For the isotropic model $K_x=K_y=K$, the self-dual condition is
\begin{equation}
\label{eq:duality-isotropic-ising-critical}
    \ee^{-2K_c}=\tanh K_c,
    \qquad
    \sinh(2K_c)=1,
    \qquad
    K_c=\frac12\log(1+\sqrt2).
\end{equation}
For the anisotropic model, the self-dual curve is
\begin{equation}
\label{eq:duality-anisotropic-ising-critical}
    \sinh(2K_x)\sinh(2K_y)=1 .
\end{equation}
This is the same condition that appeared in the transfer-matrix derivation of
the quantum--classical correspondence in the previous chapter.
In the strongly anisotropic limit, it reduces to the TFIM condition $h=J$.

\section{Operator duality in the transverse-field Ising chain}
\label{sec:operator-duality-tfim}

The quantum version of Kramers--Wannier duality is an exact change of operator
variables.  We first work on an infinite chain or in the bulk of a long open
chain, where boundary subtleties do not obscure the local algebra.  Define dual
Pauli operators on the bonds of the original lattice by
\begin{equation}
\label{eq:duality-dual-spin-definitions}
    \mu_{j+\frac12}^z
    =
    \sigma_j^x\sigma_{j+1}^x,
    \qquad
    \mu_{j+\frac12}^x
    =
    \prod_{\ell\leq j}\sigma_\ell^z .
\end{equation}
The first operator measures whether neighbouring $\sigma^x$ spins are aligned;
the second is a semi-infinite string of transverse-field operators and is called
a \emph{disorder operator}.

The operators in \cref{eq:duality-dual-spin-definitions} obey the Pauli algebra.
Indeed,
\begin{equation}
\label{eq:duality-mu-pauli-algebra}
    (\mu_{j+\frac12}^x)^2=(\mu_{j+\frac12}^z)^2=\id,
    \qquad
    \{\mu_{j+\frac12}^x,\mu_{j+\frac12}^z\}=0,
\end{equation}
and dual operators on distinct dual sites commute.  The anticommutation at the
same dual site occurs because $\mu_{j+1/2}^x$ overlaps with
$\sigma_j^x\sigma_{j+1}^x$ on exactly one site.  For different dual sites, the
number of overlaps is either zero or two, so the operators commute.

The inverse local identities are
\begin{equation}
\label{eq:duality-inverse-local-identities}
    \sigma_j^x\sigma_{j+1}^x=\mu_{j+\frac12}^z,
    \qquad
    \sigma_j^z=\mu_{j-\frac12}^x\mu_{j+\frac12}^x .
\end{equation}
Substituting these into \cref{eq:duality-chapter-tfim} gives
\begin{equation}
\label{eq:duality-tfim-dual-hamiltonian-mu}
    H_{\rm TFIM}
    =
    -J\sum_j\mu_{j+\frac12}^z
    -h\sum_j\mu_{j-\frac12}^x\mu_{j+\frac12}^x .
\end{equation}
If we rotate the dual spin axes by defining
\begin{equation}
\label{eq:duality-rotated-dual-spins}
    \widetilde{\sigma}_{j+\frac12}^x=\mu_{j+\frac12}^z,
    \qquad
    \widetilde{\sigma}_{j+\frac12}^z=\mu_{j+\frac12}^x,
\end{equation}
then the dual Hamiltonian becomes
\begin{equation}
\label{eq:duality-tfim-dual-hamiltonian-tilde}
    H_{\rm TFIM}(J,h)
    \longmapsto
    -h\sum_j
    \widetilde{\sigma}_{j-\frac12}^z
    \widetilde{\sigma}_{j+\frac12}^z
    -J\sum_j
    \widetilde{\sigma}_{j+\frac12}^x .
\end{equation}
This is the same TFIM form after relabeling the axes, with the two couplings
exchanged:
\begin{equation}
\label{eq:duality-tfim-coupling-exchange}
    (J,h)\longleftrightarrow (h,J).
\end{equation}
Equivalently, in terms of the dimensionless ratio $g=h/J$,
\begin{equation}
\label{eq:duality-g-to-inverse-g}
    g\longmapsto g^*=\frac1g,
\end{equation}
up to an overall energy scale.

\begin{remark}[Why the dual lattice is shifted]
The dual spins live on bonds $j+1/2$, not on the original sites $j$.  This shift
is not cosmetic.  It encodes the fact that the elementary dual variable is a
domain-wall variable: it naturally lives between two neighbouring order
variables.  In the continuum limit, this shift becomes invisible, but on the
lattice it is essential for getting the operator algebra and boundary conditions
right.
\end{remark}

\section{Order and disorder operators}
\label{sec:order-disorder-operators}

The duality map makes the distinction between order and disorder completely
explicit.  In the original spin language, the ferromagnetic order parameter is
$\sigma_j^x$.  In the dual language, a two-point order correlator becomes a
string of dual domain-wall variables:
\begin{equation}
\label{eq:duality-order-correlator-string}
    \sigma_i^x\sigma_j^x
    =
    \prod_{m=i}^{j-1}
    \mu_{m+\frac12}^z,
    \qquad i<j .
\end{equation}
Conversely, the dual order operator $\mu^x$ is a disorder string in the original
variables.  Its two-point function is
\begin{equation}
\label{eq:duality-disorder-correlator-string}
    \mu_{i+\frac12}^x\mu_{j+\frac12}^x
    =
    \prod_{\ell=i+1}^{j}\sigma_\ell^z,
    \qquad i<j .
\end{equation}
Thus Kramers--Wannier duality exchanges a local order operator with a nonlocal
string operator.

In the ordered phase $J>h$, the original spins have long-range $\sigma^x$ order,
while the dual variables are disordered.  In the paramagnetic phase $h>J$, the
original order is absent, but the disorder operator has long-range order:
\begin{equation}
\label{eq:duality-order-disorder-phases}
    \begin{array}{ccl}
    J>h &:& \langle \sigma^x\rangle\neq0,
    \qquad \langle \mu^x\rangle=0,\\[2pt]
    h>J &:& \langle \sigma^x\rangle=0,
    \qquad \langle \mu^x\rangle\neq0 .
    \end{array}
\end{equation}
This is the quantum-chain version of the classical statement that the low-
temperature expansion in domain walls is dual to the high-temperature expansion
in closed loops.

\section{Finite chains, constraints, and boundary conditions}
\label{sec:duality-boundary-conditions}

The local formulas above are exact in the bulk.  On a finite chain, duality also
acts on global sectors and boundary conditions.  This point is important because
many solvable-chain calculations in these notes use periodic boundary
conditions, parity sectors, and momentum quantization.

For an open chain with $L$ sites,
\begin{equation}
\label{eq:duality-open-tfim}
    H_{\rm OBC}
    =
    -J\sum_{j=1}^{L-1}\sigma_j^x\sigma_{j+1}^x
    -h\sum_{j=1}^{L}\sigma_j^z,
\end{equation}
define
\begin{equation}
\label{eq:duality-open-definitions}
    \mu_{j+\frac12}^z=\sigma_j^x\sigma_{j+1}^x,
    \qquad
    \mu_{j+\frac12}^x=\prod_{\ell=1}^{j}\sigma_\ell^z,
    \qquad
    j=1,\ldots,L-1.
\end{equation}
Then
\begin{equation}
\label{eq:duality-open-boundary-identities}
    \sigma_j^z
    =
    \mu_{j-\frac12}^x\mu_{j+\frac12}^x,
    \qquad
    \mu_{\frac12}^x\equiv\id,
    \qquad
    \mu_{L+\frac12}^x\equiv P_z,
\end{equation}
where
\begin{equation}
\label{eq:duality-global-z2-charge}
    P_z=\prod_{j=1}^{L}\sigma_j^z
\end{equation}
is the global Ising $\ZZ_2$ symmetry generator.  Therefore the dual open-chain
Hamiltonian contains boundary terms involving the fixed endpoint variables
$\mu_{1/2}^x$ and $\mu_{L+1/2}^x$.  In a fixed $P_z$ sector these are ordinary
boundary conditions for the dual chain.

For a periodic chain, the same issue appears as a relation between symmetry
sectors and boundary twists.  The definitions
\begin{equation}
\label{eq:duality-periodic-definitions}
    \mu_{j+\frac12}^z=\sigma_j^x\sigma_{j+1}^x,
    \qquad
    \mu_{j+\frac12}^x=\prod_{\ell=1}^{j}\sigma_\ell^z
\end{equation}
lead to
\begin{equation}
\label{eq:duality-periodic-constraint}
    \prod_{j=1}^{L}\mu_{j+\frac12}^z
    =
    \prod_{j=1}^{L}\sigma_j^x\sigma_{j+1}^x
    =
    \id,
\end{equation}
and
\begin{equation}
\label{eq:duality-periodic-twist}
    \mu_{L+\frac12}^x=P_z\mu_{\frac12}^x .
\end{equation}
Hence a periodic original chain maps to a dual chain with a constraint and a
boundary condition determined by the original global $\ZZ_2$ charge.  Conversely,
the global symmetry sector of the dual model determines the boundary condition
of the original model.  This is the spin-chain analogue of the parity-dependent
fermion boundary conditions encountered in the Jordan--Wigner transformation.

\begin{remark}[Duality is not always an on-site unitary]
On an infinite chain, Kramers--Wannier duality is a faithful isomorphism of local
operator algebras.  On a finite periodic chain, it is not simply an on-site
unitary map from one unconstrained $2^L$-dimensional spin Hilbert space to
another unconstrained $2^L$-dimensional spin Hilbert space.  One must keep track
of constraints, global sectors, and boundary twists.  Ignoring this point is a
common source of erroneous factors of two in partition functions and spectra.
\end{remark}

\section{Self-duality and the TFIM critical point}
\label{sec:self-duality-tfim-critical-point}

The TFIM is self-dual when $J=h$.  Since duality exchanges $J$ and $h$, the point
$J=h$ is invariant under the transformation.  The exact free-fermion spectrum of
\cref{eq:duality-chapter-tfim} is
\begin{equation}
\label{eq:duality-tfim-spectrum}
    E(k)
    =
    2\sqrt{
        (h-J\cos k)^2
        +
        (J\sin k)^2
    }.
\end{equation}
For $J,h>0$, the excitation gap is obtained at $k=0$:
\begin{equation}
\label{eq:duality-tfim-gap}
    \Delta_{\rm gap}=2\abs{h-J}.
\end{equation}
Thus the self-dual point is indeed the quantum critical point,
\begin{equation}
\label{eq:duality-tfim-critical-point}
    h=J .
\end{equation}
At this point the order and disorder descriptions are equally natural, the
Majorana mass in the continuum limit changes sign, and the scaling limit is the
Ising conformal field theory.

This conclusion can also be obtained from the two-dimensional classical Ising
model.  The transfer-matrix chapter showed that the anisotropic classical
critical condition
\begin{equation}
    \sinh(2K_x)\sinh(2K_\tau)=1
\end{equation}
reduces, in the strongly anisotropic quantum limit, to $h=J$.  Kramers--Wannier
duality and the quantum--classical correspondence therefore locate the same
critical point.

\section{When does self-duality determine a critical point?}
\label{sec:self-duality-criterion-limitations}

Self-duality is powerful, but it is not by itself a proof of criticality.  The
precise statement is conditional.

\begin{principle}[Self-duality criterion]
Suppose a model depends on a coupling ratio $g$, and an exact duality maps
$g\mapsto 1/g$.  If the model has exactly one phase transition separating two
phases exchanged by the duality, and if the transition is not hidden by an
extended critical region or by additional symmetry-breaking/topological
structure, then the transition must occur at the self-dual point $g=1$.
\end{principle}

The assumptions in this principle matter.  A self-dual point can fail to be a
unique continuous critical point for several reasons:
\begin{enumerate}[label=(\roman*)]
    \item there may be more than one transition, appearing in dual pairs away
    from the self-dual point;
    \item there may be a first-order transition at the self-dual point rather
    than a continuous critical point;
    \item there may be an extended critical phase rather than an isolated
    critical point; or
    \item the duality may exchange sectors or boundary conditions rather than
    mapping a single finite-size Hamiltonian strictly to itself.
\end{enumerate}
For the ferromagnetic TFIM with $J,h>0$, the extra spectral information in
\cref{eq:duality-tfim-spectrum} verifies that there is a single gap closing at
$J=h$.  Thus in this model self-duality is not merely suggestive; it correctly
identifies the exact critical point.

\begin{remark}[Sign conventions]
The sign of $h$ in \cref{eq:duality-chapter-tfim} is removable by the global
unitary $\prod_j\sigma_j^x$, which flips $\sigma_j^z\mapsto-\sigma_j^z$ and leaves
$\sigma_j^x\sigma_{j+1}^x$ invariant.  On a bipartite chain, the sign of $J$ can
also be changed by a staggered rotation, up to boundary subtleties for periodic
chains.  Therefore the essential ferromagnetic TFIM phase diagram is usually
drawn in terms of $\abs{h}/\abs{J}$, with the critical point at
$\abs{h}=\abs{J}$.
\end{remark}

\section{Relation to fermions and Majorana mass}
\label{sec:duality-fermions-majorana-mass}

Under the Jordan--Wigner transformation, the TFIM becomes a quadratic Majorana
or BdG Hamiltonian.  In the convention of these notes, the single-particle gap is
controlled by the difference $h-J$.  Near $k=0$ and near $h=J$, the BdG spectrum
has the relativistic form
\begin{equation}
\label{eq:duality-majorana-low-energy-spectrum}
    E(k)
    \simeq
    2\sqrt{(h-J)^2+J h\, k^2}
\end{equation}
up to higher-order lattice corrections.  The continuum theory is therefore a
massive Majorana theory with mass proportional to
\begin{equation}
\label{eq:duality-majorana-mass}
    m\propto h-J.
\end{equation}
Kramers--Wannier duality exchanges the two signs of this mass.  The ordered
Ising phase and the disordered Ising phase are therefore two massive Majorana
phases separated by the massless Majorana point $m=0$.

This observation is also the bridge to the Kitaev-chain interpretation.  The
TFIM ordered/disordered transition corresponds, after fermionization, to the
transition between the topological and trivial phases of a one-dimensional
spinless $p$-wave superconductor.  In the spin language this distinction is
expressed as order versus disorder; in the fermion language it is expressed as a
change in the Majorana mass sign and, for open chains, the presence or absence
of unpaired boundary Majorana modes.  The next chapter will formulate this same
transition in terms of winding numbers and BdG topology.

\section{Summary}
\label{sec:duality-summary}

The essential identities of this chapter are
\begin{equation}
\label{eq:duality-summary-identities}
    \mu_{j+\frac12}^z=\sigma_j^x\sigma_{j+1}^x,
    \qquad
    \mu_{j+\frac12}^x=\prod_{\ell\leq j}\sigma_\ell^z,
    \qquad
    \sigma_j^z=\mu_{j-\frac12}^x\mu_{j+\frac12}^x .
\end{equation}
They imply the bulk Hamiltonian mapping
\begin{equation}
\label{eq:duality-summary-hamiltonian-map}
    -J\sum_j\sigma_j^x\sigma_{j+1}^x
    -h\sum_j\sigma_j^z
    \quad\longleftrightarrow\quad
    -h\sum_j\widetilde{\sigma}_j^x\widetilde{\sigma}_{j+1}^x
    -J\sum_j\widetilde{\sigma}_j^z,
\end{equation}
with the two couplings exchanged.  Consequently the TFIM is self-dual at
$h=J$, and the exact spectrum confirms that this point is the unique continuous
critical point of the ferromagnetic chain.

The conceptual lesson is that duality turns nonlocal degrees of freedom into
local ones.  The same physical theory can be described either by the original
order variables $\sigma^x$, by the disorder variables $\mu^x$, by Jordan--Wigner
fermions, or by Majorana fields.  This is precisely the equivalence-map viewpoint
that underlies the organization of these lecture notes.
